\chapter{مبانی علوم شناختی محاسباتی}

\section{چارچوب مفهومی و مسئلهٔ اصلی}

\subsection{علوم شناختی چیست؟}
علوم شناختی فقط مطالعهٔ مغز به‌معنای زیستی و کالبدشناختی آن نیست. پرسش اصلی علوم شناختی این است که انسان چگونه می‌تواند \term{کنش}{Act} داشته باشد: چگونه می‌بیند، می‌شنود، حرف می‌زند، فکر می‌کند، تصمیم می‌گیرد، قضاوت می‌کند، یاد می‌گیرد، خطا می‌کند و گاهی دچار اختلال می‌شود. بنابراین موضوع علوم شناختی، فهم سازوکارهایی است که رفتار و تجربهٔ انسانی را ممکن می‌کنند.

در این نگاه، «انسان» فقط فرد بزرگسال سالم و متوسط نیست. نظریهٔ شناختی باید بتواند دربارهٔ کودک، سالمند، فرد بسیار توانمند، فرد دارای ناتوانی یا اختلال، و فردی با عملکرد عادی توضیحی منسجم بدهد. اگر نظریه‌ای فقط برای یک گروه محدود جواب بدهد، هنوز نظریهٔ عمومی خوبی دربارهٔ شناخت انسان نیست.

برای توضیح این نکته می‌توان از تعبیر «ماژول شناختی» استفاده کرد. منظور از ماژول در اینجا یک قطعهٔ سخت‌افزاری جدا و دقیقاً مرزبندی‌شده نیست؛ بلکه مجموعه‌ای از فرایندها و سازوکارهاست که باعث می‌شود انسان بتواند ادراک، زبان، حافظه، تصمیم‌گیری و عمل را انجام دهد. علوم شناختی می‌کوشد این سازوکارها را توصیف، آزمایش و در نهایت مدل کند.

\subsection{از علوم شناختی تا علوم شناختی محاسباتی}
در علوم شناختی، ممکن است یک نظریه در سطح کلامی یا مفهومی توضیح دهد که انسان چگونه یک مهارت را انجام می‌دهد. علوم شناختی محاسباتی یک گام جلوتر می‌رود و می‌پرسد آیا می‌توان آن نظریه را به شکل صوری، محاسباتی و قابل پیاده‌سازی بیان کرد یا نه.

پس هدف فقط این نیست که بگوییم «انسان از زمینه استفاده می‌کند» یا «حافظهٔ کاری محدود است». هدف این است که بتوانیم این ادعاها را به مدل، الگوریتم، معماری، آزمایش و پیش‌بینی تبدیل کنیم. یک نظریهٔ شناختی زمانی برای علوم شناختی محاسباتی مفیدتر می‌شود که بتوان آن را در قالبی اجراپذیر یا دست‌کم آزمون‌پذیر بازنویسی کرد.

به همین دلیل، بحث از نظریه‌های کلاسیک شناختی شروع می‌شود، اما در آنجا متوقف نمی‌ماند. مسیر اصلی این است که هر نظریهٔ شناختی به یک مسئلهٔ محاسباتی وصل شود: آیا این نظریه می‌تواند ضعف یک مدل هوش مصنوعی را توضیح دهد؟ آیا می‌تواند راهی برای طراحی معماری بهتر، دادهٔ بهتر یا ارزیابی بهتر پیشنهاد کند؟

\subsection{نسبت علوم شناختی محاسباتی با هوش مصنوعی}
مسیر دانشی این حوزه را می‌توان به‌صورت تقریبی چنین دید:
\[
\begin{aligned}
\text{\lr{Artificial Intelligence}}
&\rightarrow
\text{\lr{Machine Learning}}
\rightarrow
\text{\lr{Deep Learning}}\\
&\rightarrow
\text{\lr{Computational Cognitive Science}}.
\end{aligned}
\]

در هوش مصنوعی کلاسیک، با روش‌هایی مانند جست‌وجو، \eng{brute force}، \eng{BFS}، الگوریتم ژنتیک و \eng{simulated annealing} روبه‌رو می‌شویم. در یادگیری ماشین، مسئلهٔ یادگیری از داده و استخراج ویژگی برجسته می‌شود. در یادگیری عمیق، مدل‌ها بازنمایی‌های پیچیده‌تر و معماری‌های عمیق‌تر می‌سازند. در علوم شناختی محاسباتی، پرسش تازه‌ای اضافه می‌شود: کدام مهارت شناختی در این مدل‌ها غایب است و چگونه می‌توان آن را وارد مدل کرد؟

در بسیاری از دوره‌های رایج علوم شناختی محاسباتی، نمونهٔ اصلی یا \eng{case study} از حوزهٔ بینایی و پردازش تصویر انتخاب می‌شود. این انتخاب تاریخی و آموزشی قابل فهم است، چون هم علوم اعصاب و روان‌شناسی شناختی داده‌های کلاسیک زیادی دربارهٔ بینایی دارند و هم هوش مصنوعی در دوره‌ای رشد بسیار مهمی را از بینایی ماشین تجربه کرده است. با این حال، در این منبع تمرکز اصلی روی زبان، گفتار، استدلال و مدل‌های زبانی است؛ یعنی حوزه‌هایی مانند:
\begin{itemize}
  \item \term{پردازش زبان طبیعی}{Natural Language Processing / NLP}
  \item \term{استدلال}{Reasoning}
  \item \term{پردازش گفتار}{Speech Processing}
  \item \term{مدل‌های زبانی}{Language Models}
  \item \term{ترنسفورمر}{Transformer}
\end{itemize}

برای دنبال‌کردن این بحث‌ها، آشنایی پایه با \eng{NLP} مفید است، اما بسیاری از مدل‌ها و معماری‌های لازم در متن مرور می‌شوند. چون وزن اصلی مثال‌ها روی زبان و گفتار است، خواننده‌ای که از حوزه‌ای دیگر وارد می‌شود بهتر است نسبت مفاهیم مطرح‌شده را با مسئلهٔ پژوهشی یا کاربردی خود صریح کند.

\subsection{تفاوت پرسش شناختی و پرسش یادگیری ماشین}
در یادگیری ماشین، مسئله معمولاً به شکل زیر صورت‌بندی می‌شود:
\[
\text{برای این ورودی، چگونه خروجی درست را پیش‌بینی کنیم؟}
\]
در چنین چارچوبی، معیارهایی مانند \eng{accuracy}، \eng{precision}، \eng{recall} یا معیارهای مشابه اهمیت زیادی دارند. اگر مدل ورودی را بگیرد و خروجی درست بدهد، از نگاه مهندسی ممکن است موفق تلقی شود؛ حتی اگر ندانیم دقیقاً در درون مدل چه اتفاقی افتاده است.

اما در علوم شناختی، پرسش اصلی کمی متفاوت است:
\[
\text{انسان چگونه این کار را انجام می‌دهد؟}
\]
بنابراین درست بودن خروجی کافی نیست. یک مدل شناختی باید تا حدی دربارهٔ فرایند، زمان واکنش، الگوی خطا، تفاوت‌های فردی و اختلالات نیز توضیح بدهد. اگر مدلی در یک آزمون امتیاز بالا بگیرد، اما خطاهایش هیچ شباهتی به خطاهای انسانی نداشته باشد، شاید ابزار مهندسی خوبی باشد، اما الزاماً مدل شناختی خوبی نیست.

همین تفاوت باعث می‌شود علوم شناختی محاسباتی پلی میان دو دغدغه باشد: از یک طرف می‌خواهد مدل محاسباتی بسازد، و از طرف دیگر نمی‌خواهد فقط به خروجی نهایی قانع شود. مدل باید از نظر رفتاری و در صورت امکان از نظر فرایندی با شواهد انسانی قابل مقایسه باشد.

\subsection{نظریهٔ شناختی، مدل محاسباتی و اعتبارسنجی}
برای توضیح تعدد نظریه‌های علوم شناختی می‌توان از استعارهٔ فیل استفاده کرد. گویی پدیدهٔ شناخت انسان مانند فیلی بزرگ است و هر نظریه فقط بخشی از آن را لمس می‌کند: یک نظریه شنیدن را توضیح می‌دهد، نظریه‌ای دیگر حافظه را، نظریه‌ای دیگر توجه را، و نظریه‌ای دیگر تصمیم‌گیری یا خطا را. هیچ نظریه‌ای به‌تنهایی تمام ذهن انسان را توضیح نمی‌دهد.

یک نظریهٔ شناختی خوب باید چند ویژگی داشته باشد:
\begin{itemize}
  \item رفتار درست انسان را توضیح دهد.
  \item خطاهای انسان را نیز توضیح دهد، نه اینکه فقط پاسخ‌های درست را مدل کند.
  \item فقط به یک سن، یک وضعیت یا یک گروه محدود وابسته نباشد.
  \item دربارهٔ تفاوت‌های فردی و اختلالات شناختی حرفی برای گفتن داشته باشد.
  \item پیش‌بینی تولید کند و با آزمایش‌های مستقل سنجیده شود.
\end{itemize}

اعتبارسنجی در علوم شناختی معمولاً با یک مثال قانع‌کننده تمام نمی‌شود. باید آزمایش طراحی شود، دادهٔ انسانی جمع‌آوری شود، تحلیل آماری انجام شود و نظریه در موقعیت‌های مستقل دوباره سنجیده شود. اگر نظریه‌ای می‌گوید انسان در تشخیص گفتار ابتدا چند نامزد واژگانی را فعال می‌کند، باید بتواند توضیح دهد چرا در یک موقعیت نویزی یا مبهم، فرد به‌جای یک کلمه، کلمهٔ دیگری را می‌شنود.

نکتهٔ مهم این است که نظریه‌های شناختی معمولاً به شکل صفر و یکی کنار گذاشته نمی‌شوند. بسیاری از آن‌ها بخشی از پدیده را خوب توضیح می‌دهند، اما بعداً نظریه‌های کامل‌تر یا دقیق‌تر می‌آیند و بخش‌های جاافتاده را اضافه می‌کنند. بنابراین تکامل نظریه‌ها در علوم شناختی بیشتر شبیه اصلاح، گسترش و دقیق‌تر کردن است تا صرفاً رد کامل نظریهٔ قبلی.

\subsection{مغز، نواحی کارکردی و احتیاط در تفسیر}
نظریه‌های شناختی باید با شواهد مغزی نیز تا حدی سازگار باشند، اما نباید مغز را به مجموعه‌ای از جعبه‌های کاملاً جدا و بدون هم‌پوشانی تبدیل کرد. قشر مغز نواحی مختلفی دارد که با کارکردهایی مانند بینایی، شنوایی، گفتار، تعادل، هماهنگی، تصمیم‌گیری و استدلال مرتبط‌اند. برای نمونه، نواحی پس‌سری بیشتر با پردازش بینایی، نواحی گیجگاهی با بخش‌هایی از پردازش شنیداری و زبان، و نواحی پیشانی با کارکردهایی مانند برنامه‌ریزی، قضاوت و کنترل شناختی ارتباط دارند.

با این حال، این نگاشت‌ها مرزهای کاملاً تیز ندارند. کارکردهای شناختی معمولاً حاصل شبکه‌هایی از نواحی مغزی‌اند، نه فقط یک نقطهٔ منفرد. پس وقتی می‌گوییم «این ناحیه با گفتار مرتبط است» یا «این ناحیه با استدلال مرتبط است»، منظور یک رابطهٔ غالب و تقریبی است، نه مالکیت انحصاری آن ناحیه بر یک توانایی.

\subsection{کاربرد علوم شناختی و محدودهٔ این منبع}
یکی از کاربردهای مهم علوم شناختی، فهم اختلالات و بیماری‌هاست. اگر بفهمیم یک اختلال شناختی یا رشدی، مانند بخشی از طیف اوتیسم، با چه سازوکارهایی در ادراک، زبان، تعامل اجتماعی یا پیش‌بینی محیط مرتبط است، می‌توانیم مداخله‌های بهتر و انسانی‌تری طراحی کنیم. البته این منبع جایگزین آموزش تخصصی درمان اختلالات یا نوروسایکولوژی بالینی نیست.

هدف این منبع ساختن پلی میان علوم شناختی و هوش مصنوعی است. بنابراین از نظریه‌های شناختی استفاده می‌کنیم تا ضعف‌های مدل‌های محاسباتی را بهتر بفهمیم، آزمایش‌های دقیق‌تری طراحی کنیم و در صورت امکان، مدل‌هایی بسازیم که از نظر شناختی معنادارتر باشند.

\subsection{شکست مدل‌های بزرگ در مهارت‌های شناختی}
مدل‌های زبانی بزرگ در بسیاری از وظایف زبانی و دانشی عملکرد چشمگیری دارند، اما این موفقیت به‌تنهایی نشان نمی‌دهد که آن‌ها مهارت‌های شناختی را به شکل انسانی و عمومی یاد گرفته‌اند. یکی از مسیرهای مهم ارزیابی این مدل‌ها، طراحی آزمون‌هایی است که شباهت سطحی به داده‌های آموزشی را کم می‌کند و از مدل می‌خواهد رابطهٔ انتزاعی را در موقعیتی تازه تشخیص دهد.

نمونهٔ مهم، آزمون‌های \term{استدلال قیاسی}{Analogical Reasoning} با تغییرات \term{پادواقعی}{Counterfactual} است. اگر مدل در قیاس‌های آشنا خوب عمل کند، مثلاً رابطهٔ جمع بستن در «\eng{book} به \eng{books}» و «\eng{girl} به \eng{girls}» را تشخیص دهد، هنوز معلوم نیست مفهوم انتزاعی رابطه را فهمیده یا فقط الگوهای آشنا را بازتولید کرده است. برای آزمون سخت‌تر، می‌توان الفبا، نمادها، رنگ‌ها یا ترتیب‌ها را عوض کرد و دید آیا مدل هنوز رابطهٔ «بعدی بودن»، «معکوس بودن» یا «هم‌رده بودن» را منتقل می‌کند یا نه.

مثال‌هایی مانند الفبای ساختگی، داستان‌هایی با زمینه‌های متفاوت اما \eng{theme} مشترک، و جایگزینی دسته‌ها با نمادهایی مانند ضربدر یا علامت تعجب همین مسئله را نشان می‌دهند. نکتهٔ مشترک همهٔ این مثال‌ها این است که مدل نباید فقط به شباهت ظاهری تکیه کند؛ باید رابطهٔ عمیق‌تر را تشخیص دهد. پژوهش‌های جدید دربارهٔ قیاس در مدل‌های زبانی دقیقاً در همین نقطه بحث دارند: برخی نتایج افت شدید مدل‌ها را در نسخه‌های پادواقعی نشان می‌دهند، و برخی کارهای دیگر ادعا می‌کنند مدل‌های قوی‌تر یا طراحی‌های دقیق‌تر می‌توانند بخشی از این تعمیم را انجام دهند. بنابراین صورت دقیق ادعا باید محتاطانه باشد: قیاس انتزاعی و پادواقعی هنوز یکی از آزمون‌های مهم و چالش‌برانگیز برای سنجش شناختی مدل‌های زبانی است.

\subsection{تغییرات \eng{Counterfactual} و \eng{Adversarial}}
وقتی ترتیب الفبا را عوض می‌کنیم، نمادهای تازه تعریف می‌کنیم، یا یک رابطهٔ آشنا را در محیطی ناآشنا می‌گذاریم، یک تغییر \eng{counterfactual} ساخته‌ایم. اگر تغییر طوری طراحی شود که مدل را از الگوهای سطحی و حفظ‌شده دور کند، می‌تواند نقش \eng{adversarial} نیز داشته باشد؛ یعنی برای آشکار کردن ضعف مدل به کار رود.

برای انسان، تعمیم یک رابطهٔ انتزاعی به محیط تازه معمولاً طبیعی‌تر است. البته انسان هم همیشه بی‌خطا نیست، اما اگر رابطه را فهمیده باشد، می‌تواند آن را از الفبای واقعی به الفبای ساختگی منتقل کند. مدل‌هایی که بیشتر بر همبستگی‌های آماری آشنا تکیه دارند، ممکن است با همین تغییر کوچک دچار افت شدید شوند. این افت همان جایی است که پرسش شناختی آغاز می‌شود: کدام مهارت انتزاعی یا سازوکار تعمیم در مدل ضعیف است؟

\subsection{نقش دادهٔ انسانی و همبستگی خطا}
برای ارزیابی شناختی مدل‌ها، فقط نباید بپرسیم آیا مدل \eng{task} را حل کرده است یا نه. باید بپرسیم الگوی عملکرد مدل چقدر با الگوی عملکرد انسان همبستگی دارد.

دو معیار را باید از هم جدا کرد:
\begin{enumerate}
  \item آیا مدل در انجام وظیفه موفق است؟
  \item آیا سختی‌ها و خطاهای مدل با سختی‌ها و خطاهای انسان \eng{correlation} دارد؟
\end{enumerate}

این تمایز به هدف ما بستگی دارد. اگر هدف ساختن ابزاری باشد که خطاهای انسان را جبران کند، شاید نخواهیم مدل همان خطاهای انسان را تکرار کند. اما حتی در این حالت هم دانستن نسبت خطای مدل و انسان مهم است: ابتدا باید بفهمیم آیا مدل و انسان در یک خطای شناختی مشترک‌اند یا نه، سپس می‌توانیم با \eng{fine-tuning}، \eng{few-shot learning}، تغییر داده، تغییر معماری یا روش‌های مشابه آن خطا را کاهش دهیم.

اگر هدف مدل‌سازی شناخت انسان باشد، همبستگی با خطای انسانی اهمیت بیشتری پیدا می‌کند. مدلی که فقط پاسخ‌های درست بیشتری می‌دهد، اما هیچ شباهتی به الگوی خطای انسان ندارد، الزاماً توضیح‌دهندهٔ شناخت انسان نیست. به همین دلیل، در علوم شناختی محاسباتی گاهی خطاها به اندازهٔ پاسخ‌های درست مهم‌اند.

در بحث \term{هوش مصنوعی عمومی}{Artificial General Intelligence / AGI} نیز همین مسئله دیده می‌شود. برخی چارچوب‌های جدید ارزیابی، به‌جای سنجش یک مهارت واحد، مجموعه‌ای از توانایی‌های شناختی را بررسی می‌کنند. یکی از چارچوب‌های روان‌سنجی مهم برای دسته‌بندی توانایی‌های شناختی انسان، \term{نظریهٔ کتل-هورن-کرول}{Cattell-Horn-Carroll / CHC} است که توانایی‌ها را در سطوح گسترده و محدود طبقه‌بندی می‌کند. استفاده از چنین چارچوب‌هایی برای ارزیابی مدل‌های هوش مصنوعی به این معناست که مدل فقط در یک بنچمارک خاص سنجیده نشود، بلکه پروفایل شناختی آن در چند توانایی مختلف بررسی شود.

\subsection{مسیر پژوهشی}
یک مسیر پژوهشی معمول در علوم شناختی محاسباتی چنین است:
\begin{enumerate}
  \item یک ضعف یا خطای شناختی در مدل پیدا کنیم.
  \item بررسی کنیم آیا این ضعف با خطای انسان همبستگی دارد یا نه.
  \item خطای انسانی را با یک نظریهٔ شناختی توضیح دهیم.
  \item داده، آزمایش یا توصیف مناسب برای سنجش آن طراحی کنیم.
  \item مدل را با روش‌هایی مانند \eng{prompting}، \eng{few-shot learning}، \eng{fine-tuning}، \eng{RAG} یا تغییر معماری اصلاح کنیم.
  \item نشان دهیم اصلاح مدل فقط افزایش سطحی عملکرد نیست، بلکه با تحلیل شناختی مسئله ارتباط دارد.
\end{enumerate}

اگر بدون مرحلهٔ شناختی فقط مدل را بزرگ‌تر کنیم یا روی دادهٔ بیشتری آموزش دهیم، ممکن است بهبود مهندسی به دست آید، اما سهم علوم شناختی روشن نیست. ارزش پژوهشی زمانی بیشتر می‌شود که ابتدا ضعف مدل در قالب یک مهارت شناختی توضیح داده شود و سپس راهکار محاسباتی دقیقاً همان ضعف را هدف بگیرد.
